\documentclass[11pt]{article}

\usepackage[margin=1in]{geometry}
\usepackage{amsmath,amssymb,amsthm,mathtools}
\usepackage{enumitem}
\usepackage{booktabs}
\usepackage{tabularx}
\usepackage{microtype}
\usepackage[hidelinks]{hyperref}
\usepackage{xurl}

\newtheorem{definition}{Definition}
\newtheorem{lemma}{Lemma}
\newtheorem{remark}{Remark}

\title{Successor Status Envelopes and Lineage Notices for Release Maintenance:\\[0.5ex]
\large Narrow Status Tokens, Reader Pointer Order, and Public Answers}
\author{}
\date{Unified archive note v0.05 (2026-03-21; archive v489)}

\begin{document}
\maketitle

\begin{abstract}
Later papers often need one public answer after a successor release: does the successor stay on the same certified interface, and is the shipped public package complete?
Those papers should not have to choose between citing a full continuity/closure stack and citing one decorative sentence.
This note inserts one narrow bridge.
It defines a \emph{status envelope} that binds the release pair, the continuity token, the closure token, and the derived public-status token, and a \emph{lineage notice} that compresses that envelope into one reader-facing sentence plus a declared pointer order.
The result is a concise summary-stage contract: one object owns the narrow status token, another owns the outward-facing prose, and neither silently re-owns numeric deltas, role rows, or replay evidence.
\end{abstract}

\paragraph{Non-redundancy note.}
Synthesis~18 owns drift classes and continuity prerequisites.\cite{series:planstateequiv}
Synthesis~33 owns the obligation$\to$closure bridge.\cite{series:releaseclosure}
Synthesis~32 owns the stable release spine and stage order.\cite{series:releasespine}
Synthesis~31 owns the downstream support/citation/quote/clause suffix once a summary object already exists.\cite{series:downstreamnormalforms}
This note owns only the summary-stage split between narrow status tokens and outward-facing lineage notices.

\paragraph{Scope.}
This is not a new anonymity mechanism, receipt primitive, or replay protocol.
It is a publication-interface note for the small public answer emitted after continuity and package closure have already been determined.

\paragraph{Imports.}
Contract-object vocabulary is imported from Synthesis~0.\cite{series:contractobjects}
Continuity and drift language are imported from Synthesis~18.\cite{series:planstateequiv}
Closure language is imported from Synthesis~33.\cite{series:releaseclosure}
Stable stage ownership is imported from Synthesis~32.\cite{series:releasespine}
The answer-side terminal reuse / cutover discipline is imported from Synthesis~54 and Synthesis~74--75, and the concrete worked instantiation is imported from Synthesis~17.\cite{series:challengeanswerreviewmenus,series:pairedterminalcitationdiscipline,series:pairedterminalcutoverrules,series:workedexample}

\paragraph{Exports.}
This note exports (i) a successor-status envelope, (ii) a lineage-notice wrapper, (iii) a no-new-semantics compression rule for summary prose, (iv) a minimal citation boundary separating the status token from the public sentence, and (v) a local-reopen rule saying that later terse papers should reopen this note only when the exact status token or exact public sentence / pointer-order wrapper is itself live. Otherwise, once the checked-answer handoff has already been fixed, later terse papers should default to the answer-side terminal citation normal form in Synthesis~54 together with the paired terminal discipline and paired cutover rule in Synthesis~74--75.\cite{series:challengeanswerreviewmenus,series:pairedterminalcitationdiscipline,series:pairedterminalcutoverrules}

\paragraph{Later-terse-paper citation posture.}
When a later short paper only needs to answer \emph{which narrow public-status packet was declared for one concrete successor pair?}, it should stop at the minimal public status-envelope / lineage-notice card below plus the tiny status-wrapper drill. It should widen only when the live issue becomes continuity or drift classification (Synthesis~18), closure-row completion (Synthesis~33), release-stage ownership (Synthesis~32), or downstream / checked-answer terminal handoff routing (Synthesis~31, Synthesis~54, and Synthesis~74--75).\cite{series:planstateequiv,series:releaseclosure,series:releasespine,series:downstreamnormalforms,series:challengeanswerreviewmenus,series:pairedterminalcitationdiscipline,series:pairedterminalcutoverrules}

\section{Two summary-stage objects, not one}

\begin{definition}[Successor-status envelope]
For one concrete base$\to$successor comparison, a \emph{successor-status envelope}
\[
E=(\pi,\chi,\kappa,\upsilon)
\]
consists of a pair anchor
\[
\pi=(r_{\mathrm{base}},r_{\mathrm{succ}},d),
\]
a continuity token $\chi$ imported from the continuity verdict, a closure token $\kappa$ imported from the package-closure verdict, and one derived public-status token $\upsilon$ chosen for that pair.
The envelope owns the narrow public answer.
It does not own numeric diffs, role rows, or replay evidence.
\end{definition}

\begin{definition}[Lineage notice]
For one successor-status envelope $E$, a \emph{lineage notice}
\[
N(E)=(\sigma,\Pi,\eta)
\]
consists of one outward-facing status sentence $\sigma$, one ordered reader-pointer list $\Pi$ naming the first artifacts a reader should inspect if the sentence is challenged, and one reference $\eta$ back to the envelope it compresses.
The lineage notice owns presentation and navigation.
It does not replace the envelope or the narrower verdicts from which the envelope was derived.
\end{definition}

\begin{remark}[Why the split matters]
A reader-facing sentence is useful because it is short.
It is dangerous for exactly the same reason.
Without a separate envelope, later papers tend to cite one sentence as if it owned the continuity token, the package-completeness token, and the release-pair anchor all by itself.
The envelope keeps the public-status token narrow; the notice keeps the prose readable.
\end{remark}

\section{No-new-semantics compression rule}

\begin{definition}[Summary-stage compression rule]
A lineage notice $N(E)$ is \emph{summary-valid} for envelope $E$ if every contested semantic piece of its sentence $\sigma$ is implied by one field already present in $E$ or by an explicit pointer in $\Pi$ to the narrower continuity or closure owner.
In particular, a summary-valid notice may not introduce numeric deltas, role-row satisfaction claims, or replay-evidence details that are absent from the envelope and absent from the pointed owner artifacts.
\end{definition}

\begin{lemma}[Envelope-first audit rule]
Let $E$ be a successor-status envelope and let $N(E)$ be a summary-valid lineage notice.
Then any challenge to the public-status token resolves to $E$, while any challenge to the wording of the public sentence resolves first to $N(E)$ and then leftward through $\eta$ to $E$ and the pointed narrower verdicts.
\end{lemma}

\begin{proof}
By definition, the envelope owns the pair anchor, continuity token, closure token, and derived public-status token.
So any challenge to those tokens terminates at the envelope rather than at the wrapper sentence.
The lineage notice owns only the outward-facing sentence and reader-pointer order, together with a reference back to the envelope.
Because summary-validity forbids the sentence from introducing semantic content not implied by the envelope or the pointed owner artifacts, a wording challenge must either be discharged by the envelope itself or continue to the narrower continuity/closure owners named by the pointer list.
Thus the sentence stays short without becoming the owner of the underlying status semantics.
\end{proof}

\begin{remark}[Later-terse-paper citation posture]
Cite the status envelope only when the live issue is \emph{which public-status token is true} for one concrete release pair.
Cite the lineage notice only when the live issue is \emph{what public sentence readers saw} or \emph{which first pointers they were given}.
Stop at Synthesis~18 or Synthesis~33 when the live issue is the imported continuity / drift basis or the closure-ledger row / completion basis, stop at Synthesis~32 when the question has already widened to release-stage ownership, and stop at Synthesis~31 once the question has already reduced to downstream support / citation / quote / clause reuse.
Once the checked-answer handoff already imports this summary-stage basis, later terse papers should default to Synthesis~54 together with Synthesis~74--75 rather than restating the status-envelope / lineage-notice bridge locally.\cite{series:planstateequiv,series:releaseclosure,series:releasespine,series:downstreamnormalforms,series:challengeanswerreviewmenus,series:pairedterminalcitationdiscipline,series:pairedterminalcutoverrules}
\end{remark}


\paragraph{A minimal public status-envelope / lineage-notice card.}
\begin{table}[t]
\centering
\small
\begin{tabularx}{0.97\linewidth}{@{}p{0.31\linewidth}X@{}}
\toprule
\textbf{Public row} & \textbf{Pinned value}\\
\midrule
Release-pair anchor & \texttt{base\_release\_id = worked-example-draft-194}, \texttt{successor\_release\_id = worked-example-draft-194b}, and \texttt{compare\_report\_id = worked-example-compare-report-v1} \\
Imported narrow tokens & \texttt{continuity\_decision = same\_certified\_interface\_replayed\_surfaces} together with \texttt{closure\_status = complete} \\
Derived public status & \texttt{public\_status = same\_interface\_replayed\_surfaces\_complete\_package} in the worked status envelope \\
Reader-facing wrapper & The lineage notice compresses that envelope into the sentence: ``Successor remains on the same certified interface after replaying the changed fallback surfaces; required successor publication outputs are complete.'' \\
First-pointer floor & The declared reader-pointer order begins with the status envelope, compare report, successor continuity verdict, release-closure ledger, and publication-closure verdict \\
Escalation pointer & Widen to Synthesis~18 for continuity/drift basis, Synthesis~33 for row-level closure evidence, Synthesis~32 for release-stage routing, and Synthesis~31 or Synthesis~54 together with Synthesis~74--75 once the live issue is downstream reuse or checked-answer terminal handoff \\
\bottomrule
\end{tabularx}
\caption{A minimal public status-envelope / lineage-notice card. This is the smallest public answer later terse papers can quote directly when the live issue is the narrow public-status token and its declared reader-facing wrapper rather than the underlying continuity evidence, closure rows, or wider release/review routing.}
\label{tab:minimal-status-envelope-card}
\end{table}

This card is intentionally non-decorative: it pins one concrete release pair, one imported continuity token, one imported closure token, one derived public-status token, and one declared first-pointer order without silently re-owning numeric diffs, replay scope details, or row-level package-closure evidence.\cite{series:workedexample}

\paragraph{One tiny status-wrapper drill.}
In the worked Case-B successor, the narrow status packet is honest because the imported continuity token says the successor stayed on the same certified interface after replaying the changed surfaces, and the imported closure token says the required successor package is complete. That pair is exactly what the public-status token \texttt{same\_interface\_replayed\_surfaces\_complete\_package} compresses. So a later terse paper may quote the short lineage sentence and its first pointers without reopening the whole maintenance spine. But if either imported token changed --- for example if the closure token were no longer \texttt{complete} --- then the wrapper sentence would have to change here before any later paper could honestly reuse it.\cite{series:workedexample}

\begin{table}[t]
\centering
\small
\begin{tabularx}{\linewidth}{@{}l l X@{}}
\toprule
Layer & Canonical object & Question answered \\
\midrule
Narrow summary basis & status envelope & Which public-status token is true for this release pair? \\
Reader-facing wrapper & lineage notice & What short sentence and first pointers were presented to readers? \\
Narrow owners & continuity / closure objects & Which replay scope, role rows, and completeness details make that status true? \\
\bottomrule
\end{tabularx}
\caption{The summary stage stays auditable when public-status tokens and wrapper prose do not collapse into one object.}
\label{tab:statussplit}
\end{table}

\section{Worked example pointer}
The maintained worked bundle now ships both summary-stage objects.\cite{series:workedexample}
\path{example_successor_status_envelope.json} binds the base and successor release ids, the compare-report anchor, the continuity decision, the closure status, and the derived public-status token for the canned Case-B successor.
\path{example_successor_lineage_notice.json} then compresses that envelope into one public sentence and one declared reader-pointer order.
This gives later papers one narrow object to cite when they mean the status token itself, and one separate object to cite when they mean the outward-facing sentence. In the maintained worked series spine, the same status-envelope / lineage-notice pair is now imported directly onto each shipped audience/question route so later terse notes can recover the public wrapper without leaving the checked-sentence route object; once the checked-answer handoff has already been fixed, they should then default to Synthesis~54 and Synthesis~74--75 rather than restating this summary-stage bridge locally.\cite{series:challengeanswerreviewmenus,series:pairedterminalcitationdiscipline,series:pairedterminalcutoverrules}

\paragraph{Why this helps the series.}
The archive wants short papers with a stable maintenance spine.
Owning the summary-stage split here lets Synthesis~32 stay about stage order, lets Synthesis~31 stay about downstream sentence reuse, and keeps the worked example concrete instead of forcing it to act as the theory note for public-status wrappers.

\begin{thebibliography}{99}\small

\bibitem{series:contractobjects}
Anonymous.
\newblock Contract Objects for Anonymous-DHT Claims: Commitments, Menus, Receipts, and Release Bindings.
\newblock Series synthesis note (Synthesis~0), 2026. (This archive.)

\bibitem{series:planstateequiv}
Anonymous.
\newblock Plan and State Equivalence: Digest-Bound Replay Hooks and Drift Taxonomy.
\newblock Series synthesis note (Synthesis~18), 2026. (This archive.)

\bibitem{series:releaseclosure}
Anonymous.
\newblock Release Obligations and Closure Ledgers for Successor Maintenance: Required Roles, Row Satisfaction, and Completion Tokens.
\newblock Series synthesis note (Synthesis~33), 2026. (This archive.)

\bibitem{series:releasespine}
Anonymous.
\newblock Stable Release Spine Contracts for Successor Maintenance: Stage Ownership, Handoff Rules, and Auditable Walk-Throughs.
\newblock Series synthesis note (Synthesis~32), 2026. (This archive.)

\bibitem{series:challengeanswerreviewmenus}
Anonymous.
\newblock Challenge Answer Review Menus for Successor Maintenance: Smallest-Sufficient Referee Handoffs from Dockets, Profiles, Packets, and Ladders.
\newblock Series synthesis note (Synthesis~54), 2026. (This archive.)

\bibitem{series:pairedterminalcitationdiscipline}
Anonymous.
\newblock Paired Terminal Citation Discipline for Anonymous-DHT Receipts: Stable Answer Stops and Adjacent Request Reopenings.
\newblock Series synthesis note (Synthesis~74), 2026. (This archive.)

\bibitem{series:pairedterminalcutoverrules}
Anonymous.
\newblock Paired Terminal Cutover Rules for Anonymous-DHT Receipts: Stop, Widen, and Reopen as a Stable Terse-Paper Discipline.
\newblock Series synthesis note (Synthesis~75), 2026. (This archive.)

\bibitem{series:downstreamnormalforms}
Anonymous.
\newblock Downstream Normal Forms for Successor Maintenance: Summary, Support, Citation, Quote, and Clause Discipline.
\newblock Series synthesis note (Synthesis~31), 2026. (This archive.)

\bibitem{series:workedexample}
Anonymous.
\newblock Worked Example for Anonymous-DHT Receipts: Minimal Public Surface, Cross-Release Diff, and Successor Interlock.
\newblock Series synthesis note (Synthesis~17), 2026. (This archive.)

\end{thebibliography}

\end{document}
